% 可并入主论文的正文。公式、图表标签统一使用 p2- 前缀。
% 正文及独立预览默认自然分页。
\providecommand{\pTwoPageBreak}{}
\providecommand{\pTwoFigDir}{figures}
\section{问题二模型的建立与求解}
\label{sec:p2-model}

\subsection{首次观测下的干扰源可行域}

针对问题二，本文先根据首次观测确定源位置范围与第二站接收条件，
再以两站角域交集的直径评价定位效果，通过源位置采样与网格搜索确
定近似最优检测点及候选区域。

设第一个检测点为 $S_1=(x_1,y_1)$，测得示向度为 $\theta_1$。记目标圆域为
 $\Omega=\{X:\|X\|_2\le1800\}$，测角误差界为 $\alpha=1^\circ$。沿用问题
 一的角域表示，将以 $S$ 为顶点、以 $\theta$ 为中心方向的误差角域记
 为 $W(S,\theta)$。首次实测对应 $W_1=W(S_1,\theta_1)$。长度单位统一为米，距
 均取欧氏距离。

题目规定全向干扰源的有效接收半径 $R\in[1000,1500]$，且距离不超过 $5$ m 时返回近场信息而不返回示向度。既然第一站已取得示向度，源位置 $G$ 必须满足
\begin{equation}\label{eq:p2-source}
\mathcal T=W_1\cap\Omega\cap\{G:5<\|G-S_1\|_2\le1500\}.
\end{equation}
其中，角域反映测向不确定性，目标圆域排除场地外的位置，距离上限来自首次成功接收信号。数值采样用 $5$ m 作为近端极限，不把该端点视为一次有效示向度观测。

令 $u(\varphi)=(\cos\varphi,\sin\varphi)$。对于位于场地内的 $S_1$，从该点沿方向 $\varphi$ 到圆域边界的距离为
\begin{equation}\label{eq:p2-arena}
t_\Omega(\varphi)=-S_1\cdot u(\varphi)
+\sqrt{[S_1\cdot u(\varphi)]^2+1800^2-\|S_1\|_2^2}.
\end{equation}
因此，源可行域可按角度与距离参数化为
\begin{equation}\label{eq:p2-param}
G=S_1+r u(\varphi),\quad
|\varphi-\theta_1|_{\rm circ}\le\alpha,\quad
5<r\le\min\{1500,t_\Omega(\varphi)\}.
\end{equation}
圆周角差的定义与问题一一致。该区域保留了角域宽度，通常含有圆弧边界，不能直接等同于由四个端点连接成的四边形。

\subsection{第二检测点的接收约束}

令 $r_1(G)=\|G-S_1\|_2$。首次接收成功说明 $R\ge r_1(G)$，结合接收半径下限可知，给定可能源位置 $G$ 时，与已有观测相容的最小接收半径为 $\max\{1000,r_1(G)\}$。因此，保证再次接收信号的理想可行集合为
\begin{equation}\label{eq:p2-receive}
\mathcal A=\left\{S_2\in\Omega:
\|S_2-G\|_2\le\max\{1000,r_1(G)\},\quad\forall G\in\mathcal T\right\}.
\end{equation}
式\eqref{eq:p2-receive}同时考虑源位置与接收半径的不确定性；距第一站超过 $1000$ m 的源，其接收半径也必然超过 $1000$ m。

这一约束可以沿射线化简。设 $G=S_1+r u$ 且 $r>1000$，由圆域的凸性，同一射线上的 $G_0=S_1+1000u$ 也位于源可行域中。若 $\|S_2-G_0\|_2\le1000$，则
\begin{equation}\label{eq:p2-reduce}
\|S_2-G\|_2\le\|S_2-G_0\|_2+\|G_0-G\|_2
\le1000+(r-1000)=r.
\end{equation}
因此，只需约束各方向上距第一站不超过 $1000$ m 的源位置。

\pTwoPageBreak
\subsection{定位效果的评价与选址模型}

\subsubsection{由接收可行性转向定位精度}

能再次接收信号并不意味着定位区域足够小。两站观测方向接近时，角域交集可能十分狭长；适当侧向移动可改善交会几何，但效果还与观测距离有关。因此，本文沿用问题一的区域直径作为选点指标，不单独以交会角最大化代替直径最小化。图\ref{fig:p2-geometry}给出了源可行域与检测点筛选的区别。

\begin{figure}[H]
\centering
\includegraphics[width=\textwidth]{\pTwoFigDir/p2_geometry.pdf}
\caption{首次观测的源可行域与第二检测点筛选示意。左图横向放大以显示真实的 $\pm1^\circ$ 角域；右图等比例显示四点约束筛选的网格及近优候选点。}
\label{fig:p2-geometry}
\end{figure}

\subsubsection{源位置不确定时的名义交会评价}

为与现有程序一致，本文采用名义方位近似：对每个假设源位置 $G$，分别以两站指向 $G$ 的真实方位作为角域中心，同时保留 $\pm\alpha$ 的角域宽度。记 $\phi_i(G)=\operatorname{atan2}(G_y-y_i,G_x-x_i)$，则
\begin{equation}\label{eq:p2-nominal}
P_0(G;S_2)=W(S_1,\phi_1(G))\cap W(S_2,\phi_2(G)),\qquad
D_0(G;S_2)=\operatorname{diam}P_0(G;S_2).
\end{equation}
这里以首次实测确定 $\mathcal T$，再用式\eqref{eq:p2-nominal}比较不同布设的几何质量。评价角域的中心随假设源变化，故该量是名义几何指标，不等于保留首次实测结果后、对全部可能测角误差取最坏值的实际定位上界。

将源位置的不利影响纳入选址目标，得到
\begin{equation}\label{eq:p2-minmax}
J_0(S_2)=\sup_{G\in\mathcal T}D_0(G;S_2),\qquad
\inf_{S_2\in\mathcal A}J_0(S_2).
\end{equation}
对非空无界的角域交集约定 $D_0=+\infty$，并剔除对应的不良布设。与问题一保持一致，此处的定位区域仅为角域交集，不再截取场地边界，因此其直径可能超过场地直径。该指标表示相容位置间的最大分离尺度，不直接等同于某个位置估计值的实际偏差。

\pTwoPageBreak
\subsection{网格搜索与近优候选区域的构造}

\subsubsection{可行位置与源位置的离散化}

理想集合\eqref{eq:p2-receive}包含连续多个源位置约束。现有程序取角域两条边界在近端及远端的四个代表点，构造有限约束集合。对截断半径 $c$，定义
\begin{equation}\label{eq:p2-four}
V_c=\bigcup_{s=\pm1}\left\{S_1+5u(\theta_1+s\alpha),\ 
S_1+\rho_s(c)u(\theta_1+s\alpha)\right\},\quad
\rho_s(c)=\min\{c,t_\Omega(\theta_1+s\alpha)\}.
\end{equation}
用 $V_{1000}$ 中的四点检查 $\|S_2-V\|_2\le1000$。距离函数的凸性保证四点检查覆盖其凸包；但真实源区域含弧线，该凸包未必包含全部源位置，故有限约束筛选不等于完整源区域的严格探测保证。

以 $u_1=u(\theta_1)$ 为纵向单位向量，以 $n_1=(-\sin\theta_1,\cos\theta_1)$ 为侧向单位向量，写成
\begin{equation}\label{eq:p2-grid}
S_2=S_1+a u_1+b n_1,\qquad
a\in[-1800,1005],\quad b\in[-1000,1000].
\end{equation}
分别将两个区间等分为 $130$ 个网格坐标，删除场地外及不满足四点接收约束的位置，得到网格集合 $\widehat{\mathcal A}_h$。其中 $a$、$b$ 分别表示前向与侧向位移，便于解释选点结果。

误差评价不能仅检查四个端点，因为尚无依据证明最大直径一定在端点处取得。程序沿两条楔边各取 $100$ 个等距源位置，组成 $\mathcal G_N$，并计算
\begin{equation}\label{eq:p2-sampled}
\widehat J(S_2)=\max_{G\in\mathcal G_N}D_0(G;S_2),\qquad
\widehat J_{\min}=\min_{S_2\in\widehat{\mathcal A}_h}\widehat J(S_2).
\end{equation}
仅在采样评价有限的点之间比较。有限源采样可能漏掉不利位置，网格也可能错过连续最优点，故 $\widehat J_{\min}$ 只表示当前离散模型的最优值。

\subsubsection{直径计算与候选点筛选}

对每个检测点与源采样点组合，调用问题一的角域半平面表示，判断交集是否有界，再枚举边界交点并保留满足全部不等式的点。由问题一的凸性结论，区域直径等于这些可行顶点间距离的最大值。由于双站角域的边界数量固定，直接枚举点对即可，当前实现不另行调用旋转卡壳。

完整求解过程如下：
\begin{enumerate}[label=\arabic*)]
\item 根据 $S_1$、$\theta_1$ 构造接收约束代表点与源采样点。
\item 枚举局部网格并检查接收与场地约束，得到备选检测位置。
\item 遍历源采样点，以最大名义交会直径评价备选位置，剔除产生无界交集的点。
\item 比较有限评价值，输出近似最优点及评价值，并按允许偏离比例筛选候选区域。
\end{enumerate}
若备选网格点数为 $M$、源采样数为 $N$，则双站直径计算的单次规模固定，主要计算量为 $O(MN)$。取允许偏离比例 $\varepsilon\ge0$，定义
\begin{equation}\label{eq:p2-candidate}
\widehat{\mathcal C}_{\varepsilon}
=\{S_2\in\widehat{\mathcal A}_h:\widehat J(S_2)\le(1+\varepsilon)\widehat J_{\min}\}.
\end{equation}
候选面积以点数乘网格单元面积估计。$\varepsilon$ 用于放宽选点要求，不自动补偿采样误差。

\pTwoPageBreak
\subsection{算例结果与候选区域分析}

问题二未指定首次观测的具体坐标与示向度，因此采用现有程序的三个合成场景进行比较：A 为 $S_1=(0,900)$、$\theta_1=90^\circ$，B 为 $S_1=(1200,700)$、$\theta_1=120^\circ$，C 为 $S_1=(1500,900)$、$\theta_1=120^\circ$。三者分别对应沿半径朝外、斜穿圆域和近切向斜穿。所有结果均采用同一网格与源采样密度，属于名义几何算例，不是模拟器正式测试结果。

\begin{table}[H]
\centering\small
\caption{三场景的近似最优检测点与 $\varepsilon=5\%$ 候选规模}\label{tab:p2-best}
\begin{tabular}{cccccc}
\toprule
场景 & 有限评价点数 & $\widehat J_{\min}$/m & 最优 $(x_2,y_2)$/m & 候选点数 & 面积/$10^4$ m$^2$\\
\midrule
A & 3740 & 47.83 & $(565.89,1578.84)$ & 94 & 3.17\\
B & 2510 & 64.37 & $(307.83,1144.51)$ & 14 & 0.47\\
C & 3215 & 22.67 & $(1123.87,1039.86)$ & 24 & 0.81\\
\bottomrule
\end{tabular}
\end{table}

三场景通过接收约束筛选的点数分别为 $3870$、$2689$、$3318$，表\ref{tab:p2-best}仅统计其中具有有限评价值的点。A 场景关于首次观测射线对称，另一侧存在等价最优点，表中仅列一侧。坐标小数位反映网格计算结果，不表示厘米级选址精度。

\begin{figure}[H]
\centering
\includegraphics[width=\textwidth]{\pTwoFigDir/p2_candidates.pdf}
\caption{三场景的相对评价值与 $\varepsilon=5\%$ 候选区域。红线表示候选边界，星号表示一处网格最优点；颜色统一表示 $\widehat J/\widehat J_{\min}$，大于等于 $4$ 的值合并显示，空白处不参与有限值比较。}
\label{fig:p2-candidates}
\end{figure}

图\ref{fig:p2-candidates}表明，同一接收可行范围内的定位质量具有明显差异。较优位置具有一定侧向位移，可以改善沿首次方向的距离辨识能力；但前向位移与侧向位移需要联合选择，不能仅由交会角确定。场景 A 的候选区呈两侧对称分布，B、C 的近优区域则在所示条件下偏向场地内部一侧。

取 $\varepsilon=5\%$ 时，三场景候选点分别占有限评价点的 $2.51\%$、$0.56\%$、$0.75\%$，评价阈值为 $50.22$、$67.59$、$23.81$ m。这说明接收可行范围仍需进一步筛选。场景 C 的源可行距离较短，取得较小评价值；但三场景同时改变了位置和边界条件，不能据此推出普遍单调关系，也不能断言所有最优点均贴近场地边界。

\pTwoPageBreak
\subsection{参数影响与数值结果的适用范围}

\subsubsection{允许偏离比例与候选规模}

固定网格与源采样，对 $\varepsilon=2\%,5\%,10\%,20\%$ 分别筛选。由定义可知，$\varepsilon$ 增大时候选集单调扩大。图\ref{fig:p2-tradeoff}同时展示候选面积与评价阈值，反映精度要求和选点灵活性之间的取舍。
\begin{figure}[H]
\centering
\includegraphics[width=\textwidth]{\pTwoFigDir/p2_tradeoff.pdf}
\caption{允许偏离比例对候选规模与采样评价阈值的影响。折线连接四个离散计算结果，不表示已拟合连续函数关系。}
\label{fig:p2-tradeoff}
\end{figure}

以场景 A 为例，$\varepsilon$ 从 $5\%$ 增至 $10\%$ 时，面积估计由 $3.17\times10^4$ m$^2$ 增至 $7.75\times10^4$ m$^2$，评价阈值由 $50.22$ m 增至 $52.61$ m。可将 $5\%$ 作为本组算例的初始选择，再按实际可达性调整；这一比例不表示定位失败概率，也不承担测角或离散误差的保证。

\subsubsection{源位置采样的局部核验}

固定表\ref{tab:p2-best}选出的检测点，分别采用每条楔边 $50$、$100$、$200$ 点计算，并用覆盖角域内部的 $21$ 个角度、每个角度 $200$ 个距离点进行比较。表\ref{tab:p2-sampling}报告同一检测点的评价值，不对加密后的位置重新优化。
\begin{table}[H]
\centering
\caption{固定所选检测点时的采样密度比较，单位：m}\label{tab:p2-sampling}
\begin{tabular}{ccccc}
\toprule
场景 & 楔边各50点 & 楔边各100点 & 楔边各200点 & 区域内 $21\times200$ 点\\
\midrule
A & 47.7954 & 47.8261 & 48.0236 & 48.0236\\
B & 64.3667 & 64.3667 & 64.3667 & 64.3667\\
C & 22.6727 & 22.6727 & 22.6727 & 22.6727\\
\bottomrule
\end{tabular}
\end{table}

A 场景在加密后评价值增加约 $0.41\%$，说明原采样确有低估；B、C 在所列精度下保持一致。这仅支持三个所选点的局部稳定性，不能证明楔边包含全部最坏位置，也不能替代整个候选区域的网格收敛与接收约束核验。

综上，以式\eqref{eq:p2-candidate}给出当前近似模型下的第二检测点候选区域。进一步完善时，应保留首次实测角域，对第二次读数 $\phi_2(G)+e_2$ 在 $|e_2|\le\alpha$ 内取不利情形，并补充弧形源边界的接收约束核验；当前结果不作为这一完整模型的误差保证。实际返回近场信息时转入近场处理，移动与路径代价则留待问题三、四考虑。
